% Part II: Mathematical Foundations

\section{Semantic Sites}\label{sec:sites}

This section develops the mathematical foundations of Judgment Geometry from
first principles.  We construct semantic sites from programs, define the
judgment presheaf, establish the fundamental descent theorem, and prove the
key properties of the trust algebra.

\subsection{From Programs to Categories}

\begin{definition}[Semantic Site]\label{def:site}
Given a Python program $P$, the \emph{semantic site} $\Site(P) = (\mathcal{C}, J)$
consists of:
\begin{itemize}
  \item A small category $\mathcal{C}$ whose objects (\emph{coordinates})
        are the syntactic and semantic units of $P$ --- functions, classes,
        modules, scopes, variables, imports --- and whose morphisms encode
        structural relationships: restriction, inclusion, transport,
        and refinement.
  \item A Grothendieck topology $J$ on $\mathcal{C}$ specifying which families
        $\{U_i \to X\}_{i \in I}$ constitute \emph{coverings}: a class is
        covered by its methods, a module by its top-level definitions, a scope
        by the bindings it introduces.
\end{itemize}
\end{definition}

\begin{remark}
The site $\Site(P)$ is finite: a program of $n$ top-level definitions
produces $|\Ob(\mathcal{C})| = \Theta(n)$ objects and
$|\Mor(\mathcal{C})| = O(n^2)$ morphisms in the worst case (fully
connected call graph), though typical programs are sparse.
\end{remark}

The coordinate kinds implemented in \JuGeo{} mirror the Python AST:

\begin{center}
\begin{tabular}{ll}
\toprule
\textbf{CoordinateKind} & \textbf{Python Construct} \\
\midrule
\texttt{FUNCTION} & \texttt{def} / \texttt{async def} \\
\texttt{CLASS} & \texttt{class} \\
\texttt{MODULE} & Top-level file \\
\texttt{SCOPE} & Lexical scope (loop, with-block) \\
\texttt{VARIABLE} & Binding / assignment target \\
\texttt{IMPORT} & \texttt{import} / \texttt{from \ldots import} \\
\texttt{DECORATOR} & \texttt{@decorator} \\
\texttt{COMPREHENSION} & List/set/dict comprehension \\
\texttt{LAMBDA} & \texttt{lambda} expression \\
\texttt{GENERATOR} & Generator function / expression \\
\texttt{CONTEXT\_MANAGER} & \texttt{with} statement \\
\bottomrule
\end{tabular}
\end{center}

Morphisms encode four fundamental relationships:
\begin{description}
  \item[\textsc{Restriction}] $f|_U$: the behaviour of $f$ restricted to
        a sub-scope $U$.
  \item[\textsc{Inclusion}] $U \hookrightarrow V$: containment of one scope
        in another.
  \item[\textsc{Transport}] Equivalence-preserving movement of a judgment
        across an isomorphism.
  \item[\textsc{Refinement}] Strengthening a specification or narrowing a type.
\end{description}

\begin{proposition}[Category Axioms]\label{prop:category}
$\mathcal{C}$ as constructed by \JuGeo{} satisfies:
\begin{enumerate}
  \item \emph{Associativity}: $(f \circ g) \circ h = f \circ (g \circ h)$.
  \item \emph{Identity}: For each object $X$, $\mathrm{id}_X \circ f = f = f \circ \mathrm{id}_X$.
\end{enumerate}
These are verified at construction time by \texttt{CategoryStructure.check\_category\_axioms()}.
\end{proposition}

\subsection{Grothendieck Topology}\label{subsec:topology}

\begin{definition}[Covering Sieve]
A \emph{sieve} on $X \in \Ob(\mathcal{C})$ is a collection $S$ of
morphisms with codomain $X$ that is closed under precomposition:
if $f \in S$ and $g: Z \to \mathrm{dom}(f)$, then $f \circ g \in S$.
\end{definition}

\begin{definition}[Grothendieck Topology on $\Site(P)$]
The topology $J$ assigns to each $X$ a collection $J(X)$ of \emph{covering
sieves} satisfying:
\begin{enumerate}
  \item \emph{Maximality}: The maximal sieve $\{f \mid \mathrm{cod}(f) = X\}$
        is in $J(X)$.
  \item \emph{Stability}: If $S \in J(X)$ and $g: Y \to X$, then
        $g^*S = \{h \mid g \circ h \in S\} \in J(Y)$.
  \item \emph{Transitivity}: If $S \in J(X)$ and for each $f: U \to X$ in $S$
        there is $T_f \in J(U)$, then the composite sieve is in $J(X)$.
\end{enumerate}
These are verified by \texttt{SiteCoherenceChecker.check\_all()}.
\end{definition}

\subsection{The Judgment Presheaf}\label{subsec:presheaf}

\begin{definition}[Judgment 8-Tuple]\label{def:judgment}
A \emph{judgment} is a tuple
$\J = (\coord, \prop, \carrier, \evid, \oblig, \obstr, \trust, \prov)$ where:
\begin{enumerate}
  \item $\coord \in \Ob(\Site)$: the \emph{coordinate} (code element being judged),
  \item $\prop$: a \emph{proposition} about $\coord$ (postcondition, invariant, etc.),
  \item $\carrier$: the \emph{carrier type} (semantic domain of $\prop$),
  \item $\evid$: the \emph{evidence bundle} (proof artifacts, test results, solver models),
  \item $\oblig$: \emph{residual obligations} (conditions still to be discharged),
  \item $\obstr$: known \emph{obstructions} (barriers to verification),
  \item $\trust \in \Talg$: \emph{trust level} from the trust algebra,
  \item $\prov$: \emph{provenance trace} (origin and chain of custody).
\end{enumerate}
\end{definition}

The proposition kinds form five semantic categories:
\begin{description}
  \item[\textsc{Structural}] Properties of the code's shape: type annotations,
        arity, return type.
  \item[\textsc{Behavioral}] Properties of the code's execution:
        postconditions, preconditions, invariants, termination.
  \item[\textsc{Relational}] Properties relating multiple coordinates:
        equivalence, commutativity, idempotence.
  \item[\textsc{Resource}] Properties about resource usage: purity,
        effect tracking, memory safety.
  \item[\textsc{Semantic}] Higher-level properties: liveness, confluence,
        determinism.
\end{description}

\begin{definition}[Judgment Presheaf]\label{def:presheaf}
The \emph{judgment presheaf} $\mathcal{F}: \Site(P)^{\mathrm{op}} \to \mathbf{Set}$
assigns to each coordinate $\coord$ the set
$\mathcal{F}(\coord) = \{\J \mid \J.\coord = \coord\}$
of all judgments at $\coord$.  For a morphism $f: U \to V$, the restriction
$\mathcal{F}(f): \mathcal{F}(V) \to \mathcal{F}(U)$ restricts a judgment
to the sub-coordinate, possibly attenuating the trust level:
\[
  \trust(\mathcal{F}(f)(\J)) \tleq \trust(\J).
\]
\end{definition}

\begin{proposition}[Presheaf Axioms]\label{prop:presheaf}
$\mathcal{F}$ satisfies the presheaf axioms:
\begin{enumerate}
  \item $\mathcal{F}(\mathrm{id}_X) = \mathrm{id}_{\mathcal{F}(X)}$
        (identity is preserved).
  \item $\mathcal{F}(g \circ f) = \mathcal{F}(f) \circ \mathcal{F}(g)$
        (contravariant functoriality).
\end{enumerate}
\end{proposition}

\begin{proof}
Identity preservation is immediate since restricting along $\mathrm{id}$
does not change the coordinate or attenuate trust.  For functoriality,
restricting first along $g$ then $f$ produces the same judgment as
restricting along $g \circ f$, since trust attenuation composes:
$\tatten(\tatten(t, g), f) = \tatten(t, g \circ f)$.
\end{proof}

% ──────────────────────────────────────────────────────────────────────
\section{\v{C}ech Cohomology and Descent}\label{sec:descent}
% ──────────────────────────────────────────────────────────────────────

\subsection{The \v{C}ech Complex}\label{subsec:cech}

Let $\mathcal{U} = \{U_i \to X\}_{i \in I}$ be a covering family in $\Site(P)$.
The \v{C}ech complex is:
\[
  \Cech^0(\mathcal{U}, \mathcal{F}) \xrightarrow{\;\delta^0\;}
  \Cech^1(\mathcal{U}, \mathcal{F}) \xrightarrow{\;\delta^1\;}
  \Cech^2(\mathcal{U}, \mathcal{F}) \to \cdots
\]

\begin{definition}[\v{C}ech Cochain Groups]
\begin{align}
  \Cech^0(\mathcal{U}, \mathcal{F}) &= \prod_{i \in I} \mathcal{F}(U_i)
    && \text{(local sections)} \\
  \Cech^1(\mathcal{U}, \mathcal{F}) &= \prod_{i < j} \mathcal{F}(U_i \cap U_j)
    && \text{(overlap data)} \\
  \Cech^2(\mathcal{U}, \mathcal{F}) &= \prod_{i < j < k}
    \mathcal{F}(U_i \cap U_j \cap U_k)
    && \text{(triple overlaps)}
\end{align}
\end{definition}

\begin{definition}[Coboundary Map]
The coboundary $\delta^0: \Cech^0 \to \Cech^1$ sends a family of local
sections $(s_i)_{i \in I}$ to the family
$(\delta^0 s)_{ij} = s_i|_{U_i \cap U_j} - s_j|_{U_i \cap U_j}$
measuring the \emph{discrepancy} between overlapping judgments.
\end{definition}

A family $(s_i)$ is a \emph{cocycle} if $\delta^0(s_i) = 0$, meaning
the local judgments agree on all overlaps.  It is a \emph{coboundary}
if there exists a global section $s$ with $s_i = s|_{U_i}$.

\begin{theorem}[Fundamental Theorem of Judgment Geometry]\label{thm:fundamental}
Let $P$ be a program and $\prop$ a global property.  Then $P$ satisfies
$\prop$ if and only if $\Hcoh{1}(\Site(P), \mathcal{F}_\prop) = 0$.
\end{theorem}

\begin{proof}
The first \v{C}ech cohomology group is
$\Hcoh{1} = \ker(\delta^1)/\mathrm{im}(\delta^0)$.
\begin{itemize}
  \item[($\Leftarrow$)] If $\Hcoh{1} = 0$, every cocycle is a coboundary.
    Given a covering $\mathcal{U}$ and local judgments $(s_i)$ satisfying
    $\prop$ on each $U_i$ with agreement on overlaps (a cocycle), there
    exists a global section $s$ restricting to each $s_i$.  This global
    section witnesses that $P$ satisfies $\prop$.
  \item[($\Rightarrow$)] If $P$ satisfies $\prop$, choose any covering and
    restrict the global judgment to local sections.  This produces a
    coboundary, proving the cocycle is trivial.
\end{itemize}
\end{proof}

\begin{corollary}[Obstruction Interpretation]\label{cor:obstruction}
When $\Hcoh{1}(\Site(P), \mathcal{F}_\prop) \neq 0$, each non-trivial
cohomology class $[\alpha] \in \Hcoh{1}$ is an \emph{obstruction}:
a precise geometric witness to a way in which the program fails to satisfy
$\prop$.  The obstruction is localised to the overlaps where the cocycle
condition fails.
\end{corollary}

\subsection{Descent Strategies}\label{subsec:strategies}

The descent algorithm in \JuGeo{} admits four strategies, each instantiating
the abstract \v{C}ech computation differently:

\begin{description}
  \item[\textsc{Eager}] Construct local sections greedily and attempt
    gluing immediately.  Fastest for well-structured programs.
  \item[\textsc{Exhaustive}] Enumerate all minimal covering families
    and check descent for each.  Guarantees completeness.
  \item[\textsc{Iterative}] Start with a coarse cover, refine on overlap
    failures, repeat until $\Hcoh{1} = 0$ or obstruction is confirmed.
    The default strategy.
  \item[\textsc{Optimistic}] Assume gluing succeeds, verify the result
    post-hoc via a single compatibility check.  Useful when programs are
    expected to be correct.
\end{description}

\begin{theorem}[Descent Termination]\label{thm:termination}
Under the iterative strategy with depth limit $d$, descent terminates
in at most $d \cdot |\Ob(\Site)|$ refinement steps.
\end{theorem}

\begin{proof}
Each refinement strictly increases the number of patches in the covering.
Since $|\Ob(\Site)|$ is finite, the cover can be refined at most
$|\Ob(\Site)|$ times before reaching the discrete topology (every object
covers itself).  With depth limit $d$, the total step count is bounded.
\end{proof}

\subsection{The Gluing Lemma}\label{subsec:gluing}

\begin{lemma}[Gluing]\label{lem:gluing}
Let $\mathcal{U} = \{U_i \to X\}$ be a covering and $(s_i) \in \Cech^0$
a compatible family (i.e.\ a cocycle).  If $\Hcoh{1} = 0$, there exists
a unique global section $s \in \mathcal{F}(X)$ with $s|_{U_i} = s_i$ for
all $i$.  Moreover, the trust of $s$ equals
\[
  \trust(s) = \bigwedge_{i \in I} \trust(s_i) = \min_{i \in I} \trust(s_i).
\]
\end{lemma}

\begin{proof}
Existence and uniqueness follow from the vanishing of $\Hcoh{1}$.  The
trust formula follows from the conservative join axiom of the trust algebra:
the global judgment is only as trustworthy as the weakest local piece.
\end{proof}

% ──────────────────────────────────────────────────────────────────────
\section{The Trust Algebra}\label{sec:trust}
% ──────────────────────────────────────────────────────────────────────

\subsection{Definition and Axioms}

\begin{definition}[Trust Algebra]\label{def:trust}
The \emph{trust algebra} $\Talg = (T, \tjoin, \tatten, \tleq)$ is a
structure where $T$ is the linearly ordered set of trust levels:
\[
  \Tcontra \tleq \Tunver \tleq \Tcopilot \tleq \Toracle \tleq
  \Truntime \tleq \Tsolver \tleq \Tproof
\]
equipped with:
\begin{itemize}
  \item \emph{Conservative join} $\tjoin: T \times T \to T$, defined as
        $\tjoin(t_1, t_2) = \min(t_1, t_2)$.
  \item \emph{Attenuation} $\tatten: T \times \evid \to T$, satisfying
        $\tatten(t, e) \tleq t$.
  \item \emph{Promotion} $\tpromote: T \times \evid \to T$, upgrading
        trust given stronger evidence.
  \item \emph{Demotion} $\tdemote: T \to T$, downgrading trust on
        contradiction.
\end{itemize}
\end{definition}

\begin{theorem}[Trust Algebra Properties]\label{thm:trust}
The following hold in $\Talg$:
\begin{enumerate}
  \item \emph{Bounded lattice}: $(T, \tleq)$ is a bounded lattice with
        $\bot = \Tcontra$ and $\top = \Tproof$.
  \item \emph{Join properties}: $\tjoin$ is commutative, associative,
        and idempotent: $\tjoin(t, t) = t$.
  \item \emph{Attenuation monotonicity}: $t_1 \tleq t_2 \implies
        \tatten(t_1, e) \tleq \tatten(t_2, e)$.
  \item \emph{Promotion strictness}: $\tpromote(t, e) > t$ requires
        $e$ to strictly dominate the evidence at level $t$.
  \item \emph{Contradiction absorption}: $\tjoin(\Tcontra, t) = \Tcontra$
        for all $t$.
\end{enumerate}
All five properties are verified by \texttt{TrustAlgebraVerifier.verify\_all\_axioms()}.
\end{theorem}

\begin{proof}
Properties (1)--(3) and (5) follow directly from the definition of $\tjoin$
as $\min$ on a finite linear order.  Property (4) is a design invariant:
promotion requires evidence at a strictly higher tier, which is enforced
by the promotion function checking $\trust(e) > t$ before upgrading.
\end{proof}

\subsection{Trust Flow Through Morphisms}\label{subsec:trust-flow}

\begin{proposition}[Trust Attenuation Under Restriction]
For any morphism $f: U \to V$ in $\Site(P)$ and judgment $\J \in \mathcal{F}(V)$:
\[
  \trust(\mathcal{F}(f)(\J)) \tleq \trust(\J).
\]
That is, restricting a judgment to a sub-coordinate never increases its trust.
\end{proposition}

\begin{proof}
The restriction map $\mathcal{F}(f)$ applies attenuation:
the restricted judgment has less context, hence can justify no more
confidence than the original.
\end{proof}

\begin{proposition}[Global Trust]
The trust of a glued global section $s$ equals the infimum of its
local constituents:
$\trust(s) = \bigwedge_{i} \trust(s_i)$.
\end{proposition}

This ensures that \JuGeo{} never overstates verification confidence:
if any local piece has low trust, the global result inherits that caution.

\subsection{Certificate Chains}\label{subsec:certificates}

Each completed descent produces a \emph{certificate}: a cryptographic
attestation linking the program hash, the site structure, the descent
trace, and the final trust level into a tamper-evident record.

\begin{definition}[Certificate]\label{def:certificate}
A \emph{certificate} $\gamma = (h_P, h_\Site, \delta, \tau, \sigma)$
consists of:
\begin{enumerate}
  \item $h_P$: hash of the program source,
  \item $h_\Site$: hash of the site structure,
  \item $\delta$: the descent trace (covering, local sections, gluing steps),
  \item $\tau \in T$: the final trust level,
  \item $\sigma$: a digital signature over $(h_P, h_\Site, \delta, \tau)$.
\end{enumerate}
\end{definition}

Certificate chains allow downstream consumers to verify that a program
was analysed by \JuGeo{} and achieved a specific trust level, without
re-running the analysis.  This is the mechanism behind
\emph{proof-carrying Python code}.
